% The vault field guide. Build with ./build-pdf.sh
\documentclass[10pt,oneside,openany]{book}
\usepackage{vaultguide}
\InputIfFileExists{version.tex}{}{\newcommand\vaultversion{dev}\newcommand\vaultdate{\today}}

\hypersetup{pdftitle={vault: the field guide},pdfauthor={The vault Authors},
  pdfsubject={Post-quantum file encryption with a decoy passphrase}}

\begin{document}

\frontmatter
\makecover
\clearpage

% ======================================================================
\chapter*{Hi. Start here.}
\markboth{Start here}{}
\addcontentsline{toc}{chapter}{Start here}

This guide is built for skimming. Genuinely. You are allowed to jump
around, read the boxes and ignore the paragraphs, or read chapter 5 first
because that's the one with the commands in it. Nobody's checking.

Every chapter opens with a \hl{TL;DR}. Read that and you've got the gist. The
rest is there for when you want the detail, or when something goes wrong at
11pm and you need to know \emph{exactly} what a message means.

\section*{The boxes, decoded}

\begin{tldr}
The whole chapter in about three lines. If you only read one thing, read this.
\end{tldr}

\begin{quickwin}
A small thing you can do right now that makes life better.
\end{quickwin}

\begin{headsup}
Something that will trip you up if you're not expecting it.
\end{headsup}

\begin{danger}
Something that can lose you data. These are rare. Read every one of them.
\end{danger}

\begin{squirrel}
A tangent. Interesting, possibly, but you can skip it and lose nothing.
\end{squirrel}

\begin{whyso}
The reasoning behind a design decision, for the curious.
\end{whyso}

\begin{pocket}
The end-of-chapter recap. A checklist of the things worth keeping in your head.
\end{pocket}

\section*{The absolute fastest route}

Got two minutes? Here's all of it:

\begin{shell}
$ vault genkey            # make your key (pick a good passphrase)
$ vault init              # make your vault
$ vault add taxes.pdf     # put something in
$ vault list              # see what's in there
$ vault get taxes.pdf     # get it back out
\end{shell}

That's the tool. Everything after this page is detail, safety rails and the
fun decoy trick in chapter 6. \faIcon[regular]{smile}

{\hypersetup{linkcolor=vtext}\tableofcontents}

\mainmatter

% ======================================================================
\chapter{Meet vault}
\tagline{what it is, what it isn't, and a 30-second mental model}

\begin{tldr}
vault is one small program that keeps your files encrypted in a folder. It
uses post-quantum cryptography, hides your file names and sizes, and has an
optional decoy passphrase that shows a harmless set of files instead of the
real ones.
\end{tldr}

\section{What it actually does}

You give vault a file. It encrypts the file and drops the result into your
\emph{vault}, which is just a folder. Later you ask for the file back and vault
decrypts it. That's the core loop. Put in, take out.

The interesting part is what the folder looks like from the outside. Nothing.
Well, not \emph{nothing}. A pile of randomly named files, all dated 1 January
1970, all in a handful of suspiciously round sizes, none bigger than 4\,MiB,
scattered across folders named \file{00} to \file{ff}. No names. No real
sizes. No dates. None of \emph{your} folders. Big files are cut into pieces
that look like everything else. Someone who steals the folder learns roughly
how much stuff you have, and that's about it.

\begin{center}
\begin{tikzpicture}[
  doc/.style={draw=vsky,fill=vsky!8,rounded corners=2pt,minimum width=17mm,minimum height=8mm,font=\scriptsize\ttfamily},
  blob/.style={draw=vteal,fill=vteal!10,rounded corners=2pt,minimum width=24mm,minimum height=6.5mm,font=\scriptsize\ttfamily},
  lbl/.style={font=\scriptsize\labelfont,text=vslate}]
  \node[doc] (a) at (0,1.1) {taxes.pdf};
  \node[doc] (b) at (0,0) {passport.jpg};
  \node[doc] (c) at (0,-1.1) {diary.md};
  \node[lbl] at (0,2) {YOUR FILES};
  \node[circle,fill=vnavy,minimum size=19mm,text=white,font=\Large] (v) at (4,0) {\faIcon{lock}};
  \node[font=\scriptsize\displaybold,text=vnavy] at (4,-1.45) {vault};
  \node[font=\tiny,text=vslate,align=center] at (4,-1.85) {passphrase + key};
  \node[blob] (r1) at (8.2,1.3) {3fa9c1\ldots\ 1970};
  \node[blob] (r2) at (8.2,0.45) {b07e2d\ldots\ 1970};
  \node[blob] (r3) at (8.2,-0.4) {91c4aa\ldots\ 1970};
  \node[blob] (r4) at (8.2,-1.25) {e5580f\ldots\ 1970};
  \node[lbl] at (8.2,2) {WHAT A THIEF SEES};
  \foreach \n in {a,b,c} \draw[-{Stealth[length=2.5mm]},vslate,thick] (\n.east) -- (v.west);
  \foreach \n in {r1,r2,r3,r4} \draw[-{Stealth[length=2.5mm]},vteal,thick] (v.east) -- (\n.west);
\end{tikzpicture}
\end{center}

\section{The three things that make it different}

\subsection{It's post-quantum}
Most encryption you use today would, in theory, fall to a big enough quantum
computer. Nobody has one yet. But someone could copy your encrypted files
\emph{today} and sit on them until they do. It's called \emph{harvest now,
decrypt later}, and it's a real strategy.

vault uses two post-quantum algorithms at once: \hl{FrodoKEM-640} and
\hl{ML-KEM-768}. An attacker has to break \emph{both}. Chapter 9 has the gory
details.

\subsection{It hides metadata, not just contents}
Plenty of tools encrypt what's \emph{in} your files but happily leave
\file{divorce-papers-final-v3.pdf} sitting there in plain sight. vault encrypts
the names too. And the sizes. And the timestamps.

\subsection{It can lie for you}
Optionally, your key can have a second, \hl{decoy} passphrase. Type it, and
vault shows a completely separate set of files. Your real files don't appear,
don't error, don't leave a hint. They just aren't there, as far as anyone
watching can tell. That's chapter 6, and yes, it's the fun one.

\section{What vault is \emph{not}}

Honesty time. Knowing the edges matters as much as knowing the middle.

\begin{nopes}
  \item \textbf{Not a sync service.} It's a folder. You can sync the folder
        yourself (chapter 8), but vault doesn't do it for you.
  \item \textbf{Not a backup by itself.} One copy of anything is zero copies.
        Back up the vault folder \emph{and} your key.
  \item \textbf{Not magic against malware.} If something nasty is running on
        your computer while you type your passphrase, it can watch you type
        it. No file-encryption tool fixes a compromised machine.
  \item \textbf{Not recoverable without the passphrase.} There is no reset
        link. There's no ``forgot password''. That's the whole point, and it
        cuts both ways.
\end{nopes}

\begin{squirrel}
Why is it called \emph{vault} and not something cooler? Because you type it a
lot, and five letters you already know how to spell beat a clever name every
single time. Also, all the cool names were taken.
\end{squirrel}

\begin{pocket}
\begin{pocketlist}
  \item vault encrypts files into a folder and hides their names, sizes and dates.
  \item The encryption is post-quantum: two algorithms, attacker must break both.
  \item An optional decoy passphrase shows a different, harmless set of files.
  \item Lose the passphrase and the files are gone. No reset. Ever.
\end{pocketlist}
\end{pocket}

% ======================================================================
\chapter{Install it}
\tagline{grab a binary, check it, put it on your PATH}

\begin{tldr}
Download the one file for your system, check its SHA-256, rename it to
\cmd{vault} (or \cmd{vault.exe}), and put it somewhere on your \env{PATH}.
Or build every platform yourself with \cmd{./build.sh}.
\end{tldr}

\section{Pick your file}

vault is a single self-contained program. No installer, no runtime, no
dependencies. You just need the right file for your operating system and CPU.

\begin{vtable}{>{\semibold}p{34mm}Y>{\ttfamily\small}p{50mm}}
\thead{YOU HAVE} & \thead{WHICH MEANS} & \thead{DOWNLOAD} \\ \midrule
Windows PC & most laptops and desktops & vault-windows-amd64.exe \\
Windows on ARM & Surface Pro X, Snapdragon laptops & vault-windows-arm64.exe \\
Mac, 2020 or later & Apple silicon (M1, M2, M3\ldots) & vault-darwin-arm64 \\
Mac, older & Intel & vault-darwin-amd64 \\
Linux PC or server & almost everything & vault-linux-amd64 \\
Linux on ARM64 & Pi 4/5 (64-bit OS), Graviton, Ampere & vault-linux-arm64 \\
Raspberry Pi, 32-bit & Pi 2, 3, 4 on a 32-bit OS & vault-linux-armv7 \\
Raspberry Pi Zero / 1 & the tiny older ones & vault-linux-armv6 \\
Android & in Termux & vault-android-arm64 \\
FreeBSD, OpenBSD, NetBSD & amd64, arm64, 386, arm\ldots & vault-freebsd-amd64 etc. \\
Something exotic & RISC-V, POWER, s390x, MIPS, LoongArch & see \cmd{dist/} \\
\end{vtable}

\begin{squirrel}
The full release covers 43 platform builds, including AIX, Solaris, illumos,
Plan~9 and even WebAssembly. Will you ever need \file{vault-plan9-arm}?
Almost certainly not. Is it pleasing that it exists? Extremely.
\end{squirrel}

\begin{headsup}
Not sure which Linux you have? Run \cmd{uname -m}. \cmd{x86\_64} means amd64,
\cmd{aarch64} means arm64, \cmd{armv7l} means armv7, and \cmd{armv6l} means armv6.
\end{headsup}

\section{Check the file hasn't been tampered with}

Every release comes with a \file{SHA256SUMS} file. It takes ten seconds to
check, and it's the difference between running \emph{vault} and running
\emph{something pretending to be vault}. For an encryption tool, that matters.

\begin{shell}
# Linux
$ sha256sum -c SHA256SUMS --ignore-missing

# macOS
$ shasum -a 256 -c SHA256SUMS --ignore-missing

# Windows (PowerShell): compare the output with the line in SHA256SUMS
PS> Get-FileHash .\vault-windows-amd64.exe -Algorithm SHA256
\end{shell}

You want to see \cmd{OK}. Anything else: stop, and download it again from
the source you trust.

\section{Put it on your PATH}

\subsection{Linux and macOS}
\begin{shell}
$ chmod +x vault-linux-amd64
$ sudo mv vault-linux-amd64 /usr/local/bin/vault
$ vault help
\end{shell}

No admin rights? Use \file{\tld/.local/bin} instead, which most modern
distributions already have on the \env{PATH}.

\begin{headsup}
\textbf{macOS says the file ``cannot be opened''?} It's not signed by Apple,
so Gatekeeper quarantines it. Once you've checked the SHA-256, clear the flag:
\cmd{xattr -d com.apple.quarantine /usr/local/bin/vault}.
\end{headsup}

\subsection{Windows}
\begin{steps}
  \item Make a folder, say \file{C:\textbackslash Tools}, and move the file in,
        renamed to \file{vault.exe}.
  \item Start menu, search for \hl{``environment variables''}, open
        \hl{Edit the environment variables for your account}.
  \item Edit \hl{Path}, add \file{C:\textbackslash Tools}, and click OK.
  \item Open a \emph{new} terminal and run \cmd{vault help}.
\end{steps}

\subsection{Android (Termux)}
\begin{shell}
$ curl -LO https://.../vault-android-arm64
$ chmod +x vault-android-arm64 && mv vault-android-arm64 $PREFIX/bin/vault
\end{shell}

\section{Did it work?}

Run \cmd{vault help}. You should see the full command list:

\shot{help}{\cmd{vault help}: everything vault can do, on one screen}

\section{Building it yourself}

Prefer to build from source? You need Go 1.26.6 or newer. Then:

\begin{shell}
$ go build .          # just your platform, produces ./vault
$ ./build.sh          # every platform, into ./dist, with SHA256SUMS
\end{shell}

\cmd{build.sh} cross-compiles every target Go supports without C libraries.
32-bit ARM Linux gets two builds, armv6 and armv7. Want just a few? Set
\env{TARGETS}. Want them somewhere else? Set \env{OUT}.

\shot{build}{A partial build, followed by checking every binary against SHA256SUMS}

\begin{whyso}
The binaries are \emph{static} and \emph{reproducible}: no C code, no system
libraries, file paths stripped, build IDs zeroed. Two people building the same
commit get byte-identical files, so you can check a release against your own
build.
\end{whyso}

\begin{pocket}
\begin{pocketlist}
  \item One file per platform. Rename it \cmd{vault} and put it on your \env{PATH}.
  \item Always check the SHA-256 against \file{SHA256SUMS}.
  \item macOS: clear the quarantine flag after checking.
  \item \cmd{./build.sh} builds all 43 platform binaries from source.
\end{pocketlist}
\end{pocket}

% ======================================================================
\chapter{Your key}
\tagline{the one file that matters most, and the passphrase that guards it}

\begin{tldr}
\cmd{vault genkey} creates your key at \file{\tld/.vault/key}, protected by a
passphrase. It can also set a decoy passphrase. \textbf{Back the key file up.}
Without it, and its passphrase, your files cannot be recovered by anyone.
\end{tldr}

\section{Making one}

\begin{shell}
$ vault genkey
\end{shell}

vault asks for a passphrase, then asks you to type it again. Nothing appears
on screen as you type. That's deliberate, not broken.

\shot{genkey-prompt}{Waiting for your passphrase. Typing is invisible on purpose.}

Next it offers the \hl{decoy passphrase}. Press \keycap{Enter} to skip it, or
set one now. We'll cover what it's for in chapter 6. The short version: it
opens a separate, fake-but-plausible set of files.

\shot{genkey-decoy-prompt}{The optional decoy passphrase. An empty answer means ``no decoy''.}

\begin{headsup}
\textbf{The decoy can only be set here, when the key is made.} There's no
``add a decoy later'' command. If you think you might ever want one, set one
now. It costs nothing, and nobody can tell it's there.
\end{headsup}

When it's done, you get one line and one small file:

\shot{genkey}{Your key: 319 bytes, readable only by you}

\section{Choosing a passphrase that doesn't suck}

The key file is encrypted with your passphrase using \hl{Argon2id}, a
deliberately slow, memory-hungry function. Each guess costs an attacker real
time and 64\,MiB of memory. Slow guessing only helps, though, if there are a lot
of guesses to make.

\begin{checks}
  \item \textbf{Use four or more random words.} \emph{correct horse battery
        staple} beats \emph{P@ssw0rd!} by a mile, and you can actually remember it.
  \item \textbf{Make it unique.} Not your email password. Not your bank's.
  \item \textbf{Make the decoy believable.} It should feel like something
        you'd genuinely pick. More in chapter 6.
\end{checks}

\begin{quickwin}
Use a password manager's \emph{passphrase} generator (most have one) to make
the words, then practise typing it five times. Muscle memory is a legitimate
security tool.
\end{quickwin}

\section{Where it lives}

By default your key is \file{\tld/.vault/key}. On Windows, that's
\file{\%USERPROFILE\%\textbackslash.vault\textbackslash key}. Want it
somewhere else, like a USB stick? Two options:

\begin{shell}
$ vault genkey --key /media/usb/vault.key     # this once
$ export VAULT_KEY=/media/usb/vault.key       # every command from now on
\end{shell}

\begin{squirrel}
Keeping the key on a USB stick and the vault on your laptop is a neat trick.
Pull the stick out and the vault is inert, even if someone knows your
passphrase. Two things to steal instead of one.
\end{squirrel}

\section{Back. It. Up.}

\begin{danger}
\textbf{If you lose the key file, or forget the passphrase, your files are
gone.} Not ``hard to get back''. Gone. vault has no recovery mechanism, no
master key, no back door. That's what makes it trustworthy, and it's also
what makes a lost key final.
\end{danger}

The key file is encrypted, so a copy is safe to keep somewhere a little less
private than your laptop. Some decent options:

\begin{itemize}
  \item A USB stick in a drawer. Or two, in two different drawers.
  \item Your password manager, as an attached file.
  \item Printed out. It's only a few lines of text. Paper survives surprisingly well.
\end{itemize}

\begin{whyso}
\textbf{What's actually in the key file?} Two equal-sized encrypted slots, in
random order. Each passphrase opens exactly one. When there's no decoy, the
spare slot is just random bytes, and nobody can tell random bytes from an
encrypted slot. So the file gives away nothing about whether a decoy exists.
\end{whyso}

\begin{pocket}
\begin{pocketlist}
  \item \cmd{vault genkey} makes \file{\tld/.vault/key}. It won't overwrite an existing key.
  \item Four-plus random words. Unique. Practise it.
  \item Decide on a decoy now. It can't be added later.
  \item Back up the key file. Somewhere else. Today.
\end{pocketlist}
\end{pocket}

% ======================================================================
\chapter{Your first vault}
\tagline{one command, and a peek inside}

\begin{tldr}
\cmd{vault init} creates an empty vault at \file{\tld/.vault/store}. It's
just a folder. Inside it seeds some random filler, which is normal and part of
the decoy trick.
\end{tldr}

\section{Create it}

\begin{shell}
$ vault init
\end{shell}

That's it. No passphrase needed, because an empty vault has nothing to
protect yet.

\shot{init}{A fresh vault. Two sub-folders, a marker, and\ldots\ records already?!}

Inside \file{records/} and \file{objects/} are 256 folders each, named
\file{00} to \file{ff}. Every encrypted file lives in the one matching the
first two characters of its name. That just keeps any one folder from
getting huge, which some file systems and sync services dislike. All 256
exist from the start, so the layout itself gives nothing away.

\section{Wait, why is there already stuff in it?}

Good catch. A brand new vault holds somewhere between 8 and 32 entries of
random filler. They're encrypted to a key that was thrown away the moment
it was used, so \emph{nobody} can open them. Not you, not us, not anyone.

It looks weird, but it's what makes the decoy passphrase believable. Every
vault, decoy or not, always contains entries you can't open. So when someone
with your decoy passphrase sees unopenable entries, that's completely
unremarkable. It's just filler. (Some of it might not be. They can't tell.)

\begin{headsup}
By default the filler takes a few tens of megabytes. That's the price of plausible
deniability, and it's paid once.
\end{headsup}

\section{More filler, for bigger secrets}

A few tens of megabytes of filler can only cover for a few files. If you
plan to hide more than that behind a decoy, ask for more when you create the
vault:

\begin{shell}
$ vault init --filler 2G
\end{shell}

vault keeps writing filler until there's at least that much. Sizes take
\cmd{K}, \cmd{M}, \cmd{G} or \cmd{T}. The filler entries are built exactly
like real ones, from a mix of small and large random ``files'', so a big
file of yours sits among big files of filler. It costs the disk space and a
few seconds per gigabyte, once. It can only be chosen at \cmd{init}: added
later, the filler would show up as a burst of new entries, which is exactly
what it's meant to hide.

\begin{quickwin}
Rule of thumb: set \cmd{--filler} to at least as much as you ever expect to
hide from the decoy passphrase. Without a decoy, you don't need it.
\end{quickwin}

\section{More than one vault}

Want a separate vault for work stuff? Point \cmd{--vault} anywhere:

\begin{shell}
$ vault init --vault ~/work-vault
$ vault add contract.pdf --vault ~/work-vault
\end{shell}

Or set \env{VAULT\_DIR} so every command uses it. The same key works with any
number of vaults.

\begin{quickwin}
Add \cmd{export VAULT\_DIR=\ldots} to your shell profile and you'll never
type \cmd{--vault} again.
\end{quickwin}

\begin{pocket}
\begin{pocketlist}
  \item \cmd{vault init} makes \file{\tld/.vault/store}. \cmd{--vault DIR} to put it elsewhere.
  \item Filler entries are normal. Nobody can open them. That's the point.
  \item Hiding a lot behind a decoy? \cmd{vault init --filler 2G} (or however much you'll hide).
  \item One key, as many vaults as you like.
\end{pocketlist}
\end{pocket}

% ======================================================================
\chapter{Everyday moves}
\tagline{add, list, get, mv, rm: the five commands you'll actually use}

\begin{tldr}
\cmd{add} puts files in (optionally into a folder), \cmd{list} shows them,
\cmd{get} takes one back out, \cmd{mv} renames, \cmd{rm} deletes. Each asks
for your passphrase. That's the whole daily workflow.
\end{tldr}

\section{Putting things in: \texttt{add}}

\begin{shell}
$ vault add FILE-OR-FOLDER [FOLDER-IN-VAULT]
\end{shell}

Give it a file, and it goes in under its own name. Give it a second argument,
and it goes into that folder inside the vault. Give it a whole folder, and
everything inside comes along too, subfolders and all.

\shot{add}{Adding a file, a file into a folder, a whole folder, and another file}

Notice what happened there:

\begin{itemize}
  \item \file{Documents/tax-return-2025.pdf} became plain \file{tax-return-2025.pdf}.
        vault keeps the \emph{file name}, not the path you typed.
  \item With a second argument, \file{id}, the passport scan landed in
        \file{id/passport-scan.jpg}.
  \item Adding the folder \file{Documents/letters} kept its own name and
        structure: \file{letters/bank.txt}, \file{letters/landlord.txt}.
\end{itemize}

\begin{headsup}
\textbf{Folders inside a vault aren't real folders.} They're just part of each
file's name, like \file{letters/bank.txt}. So there's no \cmd{mkdir}, and a
folder disappears when its last file goes. No housekeeping required.
\end{headsup}

\subsection{Adding something that's already there}
Adding a name that already exists \hl{replaces} it, and vault tells you so.
The old version is deleted \emph{after} the new one is safely stored, so a
crash halfway through never leaves you with neither.

\subsection{What gets skipped}
When adding a folder, anything that isn't a plain file, such as symbolic
links, sockets or devices, is skipped with a note. Files larger than 1\,TiB are
refused. (If you have a single 1\,TiB file, we have questions.)

\section{Seeing what's inside: \texttt{list}}

\begin{shell}
$ vault list            # everything
$ vault list letters    # just what's inside letters/
\end{shell}

Each line shows the size, the file's \emph{own} last-modified time from when
you added it, and its name, sorted alphabetically.

\shot{list}{The full catalogue, then just one folder}

\begin{whyso}
\textbf{Why does \texttt{list} need my passphrase?} Because the names are
encrypted too. Without the key, vault can't read the catalogue any better
than a thief can.
\end{whyso}

\section{Getting things out: \texttt{get}}

\begin{shell}
$ vault get NAME               # writes the file here, under its own name
$ vault get NAME -o OUTPUT     # writes it wherever you say
\end{shell}

\shot{get}{Getting files back, including vault refusing to overwrite}

Three safety habits are built in, so you don't need to think about them:

\begin{checks}
  \item \textbf{It never overwrites.} If the output exists, vault stops. Move
        the old one or use \cmd{-o}.
  \item \textbf{Nothing appears until everything checks out.} vault writes to a
        hidden temporary file, verifies every single byte, and only
        then gives it its real name. A tampered or damaged file never shows up
        half-written.
  \item \textbf{It restores the details.} You get back the original permissions
        and modification time.
\end{checks}

\section{Renaming and moving: \texttt{mv}}

\begin{shell}
$ vault mv OLD NEW
\end{shell}

\cmd{NEW} is the \emph{full} new name, folder included. So to move
\file{notes.md} into \file{ideas}, it's \cmd{vault mv notes.md
ideas/notes.md}, not \cmd{\ldots\ ideas/}. It won't overwrite an existing name,
and it doesn't re-encrypt the file, so it's instant however big the file is.

\section{Deleting: \texttt{rm}}

\begin{shell}
$ vault rm NAME
\end{shell}

Removes the file's catalogue entry and its encrypted contents. One file per
command, on purpose. There's no recursive delete to fat-finger.

\shot{mv-rm}{Replace, move, delete, and the result}

\begin{danger}
\textbf{\texttt{rm} is immediate and there's no undo.} No recycle bin, no
``are you sure?''. If it matters, make sure it's backed up somewhere first.
\end{danger}

\begin{squirrel}
Need to delete a whole folder? Loop over \cmd{list}. The name starts at
character 31 of each line, so:\\
\cmd{vault list old | cut -c31- | while read -r f; do vault rm "\$f"; done}\\
With \env{VAULT\_PASSPHRASE} set, or it'll ask every time. See chapter 7.
\end{squirrel}

\section{Flags work anywhere}

Every command understands \cmd{--key PATH} and \cmd{--vault DIR}, and they can
go before, between or after the other arguments. These all do the same thing:

\begin{shell}
$ vault add --vault ~/work report.pdf
$ vault add report.pdf --vault ~/work
$ vault add report.pdf --vault=~/work
\end{shell}

\begin{pocket}
\begin{pocketlist}
  \item \cmd{add FILE [FOLDER]}: keeps the file's name, optionally in a folder. Folders add recursively.
  \item Adding an existing name replaces it, safely.
  \item \cmd{list [FOLDER]}, \cmd{get NAME [-o OUT]}, \cmd{mv OLD NEW}, \cmd{rm NAME}.
  \item \cmd{get} never overwrites and never publishes a half-checked file.
  \item \cmd{rm} has no undo.
\end{pocketlist}
\end{pocket}

% ======================================================================
\chapter{The decoy passphrase}
\tagline{one key, two passphrases, two completely different vaults}

\begin{tldr}
Set a decoy passphrase when you run \cmd{vault genkey}. Every command run with
the decoy sees only decoy files. Your real files are invisible to it: they
don't list, can't be fetched, and look exactly like the random filler every
vault has. Make the decoy believable, and keep it at least as big as the real
thing.
\end{tldr}

\section{The idea}

Sometimes ``I won't tell you my passphrase'' isn't an option. A border
crossing. A pushy partner. Someone standing over you. The decoy passphrase is
for that moment. You type it, the vault opens, and it's full of genuinely
private-looking stuff: bank letters, some personal notes, a few scans.
Everything behaves normally. It's just not your \emph{real} stuff.

\section{How it looks}

Here's one vault, opened twice, with the two different passphrases:

\shot[1.2]{decoy-views}{Same key file, same vault folder, two passphrases, two worlds}

Nothing on the right hints that the left exists. Ask the decoy for a real
file and you get the same message you'd get for a typo. The folder's entry
count doesn't give it away either, because filler is always there:

\shot{decoy-hidden}{The decoy can't see real files, and a few dozen entries on disk is unremarkable}

\begin{center}
\begin{tikzpicture}[x=5.4mm,y=5.4mm]
  \def\cells{FRDFFRDRFFDRFDFRRFDFFRDFRDFFRFDRFDFFRDRF}
  \newcommand\grid[4]{% x offset, title, colour for R, colour for D
    \node[font=\scriptsize\labelfont,text=vslate,anchor=south] at (#1+3.5,5.35) {#2};
    \foreach \i in {0,...,39}{
      \pgfmathtruncatemacro\cx{mod(\i,8)}\pgfmathtruncatemacro\ry{\i/8}
      \StrChar{\cells}{\numexpr\i+1\relax}[\kind]
      \if\kind R\def\col{#3}\else\if\kind D\def\col{#4}\else\def\col{vline}\fi\fi
      \fill[\col,rounded corners=1pt] (#1+\cx,4-\ry) rectangle ++(0.82,0.82);}}
  \grid{0}{ON DISK}{vline}{vline}
  \grid{10}{REAL PASSPHRASE}{vteal}{vline}
  \grid{20}{DECOY PASSPHRASE}{vline}{vviolet}
\end{tikzpicture}\\[2pt]
{\footnotesize\color{vslate}Every square is an encrypted entry. Each passphrase lights up only its own. Grey is filler, or the other side.}
\end{center}

\section{Setting it up}

The decoy is chosen when the key is made, and only then:

\begin{steps}
  \item Run \cmd{vault genkey}. Type your \emph{real} passphrase, twice.
  \item At \hl{Enter decoy passphrase (empty for none)}, type the decoy. Twice.
  \item Run \cmd{vault init}, adding \cmd{--filler} with at least as much as
        you plan to hide, say \cmd{--filler 5G}.
  \item Add some decoy files, typing the \emph{decoy} passphrase when asked.
  \item Add your real files, typing the \emph{real} passphrase.
\end{steps}

The two must be different (vault checks). In scripts, set
\env{VAULT\_DECOY\_PASSPHRASE} alongside \env{VAULT\_PASSPHRASE} before
\cmd{genkey}. Chapter 7 has more on scripting.

\begin{headsup}
\textbf{Already have a key without a decoy?} You'll need a new key: run
\cmd{genkey} with a different \cmd{--key}, make a new vault, \cmd{get}
everything out of the old one and \cmd{add} it to the new. The key file
format deliberately gives no way to tell whether a decoy slot is in use, so
there's nothing to safely ``upgrade''.
\end{headsup}

\section{Why it holds up}

\begin{itemize}
  \item \textbf{The key file can't tell.} It always has two slots of identical
        size in random order. With no decoy, one slot is random noise, which
        looks the same as an encrypted slot. Both are tried every time, so
        even the timing is the same.
  \item \textbf{The vault can't tell.} Real and decoy entries are sealed to
        different keys, but they're the same size, the same shape, and all
        dated 1970. Nothing on disk says which key an entry belongs to.
  \item \textbf{The filler covers for you.} Every vault starts with at
        least 8 to 32 entries nobody can open, and as much more as you ask
        for with \cmd{--filler}. So ``some entries I can't open'' is true of
        \emph{every} vault, decoy or not.
  \item \textbf{The errors are identical.} A real file asked for by name with
        the decoy passphrase gives exactly the same ``not in the vault'' as
        a typo.
\end{itemize}

\section{Making it believable: the playbook}

A decoy only works if it looks like the real thing. An empty decoy vault, or
one with two files called \file{test.txt}, is a neon sign.

\begin{checks}
  \item \textbf{Fill it with plausibly private stuff.} Bank letters, payslips,
        a diary-ish notes file, scans of documents. The kind of thing someone
        \emph{would} encrypt.
  \item \textbf{Keep it bigger than your real collection}, or at least as big.
        See the next section for why.
  \item \textbf{Keep it alive.} Add to it now and then, so the dates in
        \cmd{list} look like a real person uses it.
  \item \textbf{Know it cold.} Practise the decoy as much as the real one.
        Hesitating on the ``real'' passphrase is a tell.
  \item \textbf{Never mention it.} Not in notes, not in a README, not
        in a script sitting on the same machine.
\end{checks}

\section{The honest limits}

\begin{headsup}
\textbf{Numbers can give you away.} Someone holding the decoy passphrase
can count the records and pieces they can't open. Without \cmd{--filler}
there are only a few tens of megabytes of filler, so if \file{records/} holds
400 entries, or \file{objects/} holds a 300\,MB video's 72 pieces, the
filler can't explain the difference. They can't see \emph{what} the
leftovers are, but they can see there's more than filler. Create the vault
with \cmd{--filler} at least as large as what you'll hide, and keep the
decoy collection believable next to it.
\end{headsup}

\begin{headsup}
\textbf{Growth over time is visible.} If someone copies your vault folder
today and again next month, they can see it grew, even if the decoy didn't
change. They can't see \emph{what} was added. They can see \emph{that}
something was.
\end{headsup}

\begin{nopes}
  \item It won't help if someone saw you type the real passphrase, or has
        malware on your machine.
  \item It doesn't stop anyone deleting the raw files. It hides; it doesn't guard.
  \item It isn't legal advice. Some places treat misleading an official as an
        offence in its own right. Know where you are.
\end{nopes}

\begin{pocket}
\begin{pocketlist}
  \item The decoy is set at \cmd{genkey} time only.
  \item Decoy and real are two separate worlds in one folder. Neither can see the other.
  \item Filler, identical entry sizes and a two-slot key file make the hiding work.
  \item Keep the decoy big, believable and alive. Keep quiet about it.
\end{pocketlist}
\end{pocket}

% ======================================================================
\chapter{Scripts and automation}
\tagline{environment variables, exit codes, and a few recipes}

\begin{tldr}
Set \env{VAULT\_PASSPHRASE}, \env{VAULT\_KEY} and \env{VAULT\_DIR} and vault
never asks anything. Output you'd pipe (\cmd{list}, \cmd{help}) goes to
standard output; everything else goes to standard error. Exit code 0 means
it worked.
\end{tldr}

\section{The environment variables}

\begin{vtable}{>{\leavevmode\ttfamily\small\color{vvioletdeep}}p{50mm}Y}
\thead{VARIABLE} & \thead{WHAT IT DOES} \\ \midrule
VAULT\_PASSPHRASE & The passphrase. Skips the prompt. With \cmd{genkey}, it's the new key's passphrase. \\
VAULT\_DECOY\_PASSPHRASE & \cmd{genkey} only: the decoy passphrase. If \env{VAULT\_PASSPHRASE} is set and this isn't, no decoy is made. \\
VAULT\_KEY & The key file. Default \file{\tld/.vault/key}. Same as \cmd{--key}. \\
VAULT\_DIR & The vault folder. Default \file{\tld/.vault/store}. Same as \cmd{--vault}. \\
\end{vtable}

\shot{scripting}{With the variable set, no prompts. A failed command exits with status 1.}

\begin{danger}
\textbf{Don't type a passphrase straight into a command line.} It lands in your
shell history. Read it without echo instead:\\[2pt]
\cmd{read -rs VAULT\_PASSPHRASE \&\& export VAULT\_PASSPHRASE}\\[2pt]
And remember that other programs running as you can read your environment.
For everyday use, just let vault prompt you.
\end{danger}

\section{Exit codes and output}

\begin{vtable}{>{\ttfamily\small}p{16mm}Y}
\thead{CODE} & \thead{MEANING} \\ \midrule
0 & It worked. \\
1 & Something went wrong. The reason is on standard error, starting \cmd{vault:}. \\
2 & You ran \cmd{vault} with no command at all. The usage text is on standard error. \\
\end{vtable}

\cmd{list} and \cmd{help} write to standard output. Progress lines like
\cmd{added \ldots}, \cmd{replaced \ldots} and every error go to standard
error. So \cmd{vault list > files.txt} gives you a clean file list.

\begin{squirrel}
The \cmd{list} columns are fixed-width: size in 10 characters, two spaces,
the date and time in 16, two spaces, then the name. The name therefore always
starts at character 31, which is why \cmd{cut -c31-} works even for names with
spaces in.
\end{squirrel}

\section{Recipes}

\subsection{Nightly: stash a folder}
\begin{shell}
#!/bin/sh
# ~/bin/stash-documents: re-adds everything in ~/Documents/private
export VAULT_PASSPHRASE="$(cat ~/.config/vault-pass)"   # chmod 600 this file!
vault add ~/Documents/private archive || exit 1
\end{shell}

Then run it from cron, say at 2am: \cmd{0 2 * * * \$HOME/bin/stash-documents}.
Re-adding replaces each file, so it always holds the latest copy.

\begin{headsup}
Keeping the passphrase in a file trades safety for convenience. Anyone who can
read that file and the key can open the vault. Only do it on a machine you
trust completely, and never for the decoy.
\end{headsup}

\subsection{Get everything back out}
\begin{shell}
#!/bin/sh
# restore the whole vault into ./restore, keeping its folders
vault list | cut -c31- | while IFS= read -r name; do
  mkdir -p "restore/$(dirname "$name")"
  vault get "$name" -o "restore/$name"
done
\end{shell}

\subsection{PowerShell on Windows}
\begin{shell}
PS> $env:VAULT_DIR = "D:\vault"
PS> vault add .\report.docx
PS> vault list
\end{shell}

\begin{pocket}
\begin{pocketlist}
  \item \env{VAULT\_PASSPHRASE}, \env{VAULT\_KEY}, \env{VAULT\_DIR}: no prompts, no flags.
  \item Exit 0 is success, 1 is an error, 2 is ``no command''.
  \item \cmd{list} goes to stdout, so it pipes cleanly. Names start at column 31.
  \item Keep passphrases out of your shell history.
\end{pocketlist}
\end{pocket}

% ======================================================================
\chapter{Locking, syncing and backups}
\tagline{two terminals, two machines, and keeping copies}

\begin{tldr}
vault locks itself: one writer at a time, and readers wait for writers. You
don't have to do anything. Syncing the folder between machines works if only
one machine changes it at a time. Back up the folder and the key, and you're
covered.
\end{tldr}

\section{Locking is automatic}

Commands that change the vault (\cmd{add}, \cmd{mv}, \cmd{rm}) take an
\hl{exclusive} lock. Commands that only read (\cmd{list}, \cmd{get}) take a
\hl{shared} one, so several can read at once. If a command has to wait, it
says so, then carries on by itself:

\shot[1.2]{locking}{The second writer waits politely for the first to finish}

The lock is taken \emph{after} you've typed your passphrase, so a colleague
who's slow at typing never blocks you.

\begin{whyso}
\textbf{What if vault crashes while holding the lock?} On Linux, macOS, the
BSDs and Windows, the operating system drops the lock the moment the process
dies. Nothing to clean up. On Plan~9 and WebAssembly, which have no file
locks, vault uses a \file{.lock} file instead. If one's ever left behind,
vault tells you after a minute exactly which file to delete.
\end{whyso}

\section{Syncing between machines}

A vault is just a folder, so Syncthing, Dropbox, iCloud Drive, OneDrive,
\cmd{rsync} or a USB stick will all carry it between computers. Copy your key
file to each machine too, over something private.

\begin{headsup}
\textbf{Locks don't travel through sync services.} Two machines writing at
the same time won't wait for each other. The rule: \hl{one machine writes at a
time, and let the sync finish before you write on the other one.}
\end{headsup}

If you do slip up, it's rarely a disaster:

\begin{itemize}
  \item Sync ``conflicted copy'' files have names vault doesn't recognise, so
        it simply ignores them.
  \item If both machines added the same name, \cmd{list} shows the newest,
        and the next \cmd{add} or \cmd{rm} of that name tidies up the other.
\end{itemize}

Network file systems are fine too. NFS honours vault's locks on Linux and the
BSDs, and Windows file shares honour them from Windows.

\section{Backups}

Two things to back up: the \hl{vault folder} and the \hl{key file}. Ideally
in different places. Any backup tool works, because it's all ordinary files.

\begin{quickwin}
vault never edits a file in place. Every entry is written once and later
deleted, never changed. That makes incremental backups (\cmd{rsync},
restic, borg, Time Machine) fast and tidy: each run only copies what's new.
\end{quickwin}

\begin{shell}
$ rsync -a --delete ~/.vault/store/ /media/backup/vault-store/
\end{shell}

\begin{headsup}
A backup of the vault is only as good as your ability to open it. Test it:
point \cmd{--vault} at the backup copy and run \cmd{vault list} now and then.
\end{headsup}

\begin{pocket}
\begin{pocketlist}
  \item Locking is automatic. Writers take turns; readers share.
  \item Sync across machines freely, but write on one machine at a time.
  \item Back up the folder and the key, separately. Test the backup.
\end{pocketlist}
\end{pocket}

% ======================================================================
\chapter{Under the hood}
\tagline{the cryptography, the file layout, and exactly what a thief learns}

\begin{tldr}
Each file gets its own random 256-bit key, wrapped with a hybrid of FrodoKEM
and ML-KEM, and is encrypted in authenticated 64\,KiB chunks. Catalogue
entries are padded to one fixed size, file contents are cut into 4\,MiB
pieces with the last padded to a power-of-two bucket, and every timestamp on
disk is set to 1970. An attacker learns roughly how
much you store. They can delete or roll back entries, but not read or forge them.
\end{tldr}

\section{The cryptography}

\begin{vtable}{>{\semibold}p{38mm}Y}
\thead{JOB} & \thead{HOW} \\ \midrule
Protecting your key & Argon2id (3 passes, 64\,MiB, 4 lanes) turns the passphrase into a key that seals your identity with XChaCha20-Poly1305. \\
Your identity & A 32-byte seed from which the FrodoKEM-640 and ML-KEM-768 key pairs are derived. \\
Per-file key & 256 random bits, fresh for every single entry. \\
Wrapping that key & FrodoKEM-640 \emph{and} ML-KEM-768 encapsulations, combined with HKDF-SHA-256 into one key-encryption key, then ChaCha20-Poly1305. \\
Header integrity & HMAC-SHA-256 under the per-file key. \\
The contents & ChaCha20-Poly1305 in a STREAM construction: 64\,KiB chunks, each authenticated, with the last one marked so truncation is caught. \\
\end{vtable}

\begin{whyso}
\textbf{Why two post-quantum algorithms?} ML-KEM is the new standard,
efficient and heavily analysed. FrodoKEM is slower and bulkier, but rests on
older, more conservative maths. Combining them means a breakthrough against
either one alone isn't enough.
\end{whyso}

\section{What's in the folder}

\begin{center}
\begin{tikzpicture}[
  dir/.style={font=\small\ttfamily\color{vskydeep},anchor=west},
  f/.style={font=\small\ttfamily\color{vtext},anchor=west},
  note/.style={font=\footnotesize\color{vslate},anchor=west},
  tree/.style={draw=vslate!45,line width=0.8pt}]
  \node[dir] (root) at (0,0) {\tld/.vault/store/};
  \node[f] (m) at (0.6,-0.7) {.vault};      \node[note] at (4.2,-0.7) {marker: ``VAULT 2''};
  \node[f] (l) at (0.6,-1.3) {.lock};       \node[note] at (4.2,-1.3) {lock file, always empty};
  \node[dir] (r) at (0.6,-1.9) {records/0f/};  \node[note] at (4.2,-1.9) {one encrypted catalogue entry per file};
  \node[f] (r1) at (1.2,-2.5) {0f1170a6\ldots}; \node[note] at (4.2,-2.5) {every one exactly 18,704 bytes};
  \node[dir] (o) at (0.6,-3.1) {objects/3e/};  \node[note] at (4.2,-3.1) {file contents, in 4\,MiB pieces};
  \node[f] (o1) at (1.2,-3.7) {3e78275e\ldots}; \node[note] at (4.2,-3.7) {64\,KiB, 128\,KiB\ldots\ 4\,MiB, never in between};
  \draw[tree] (0.3,-0.25) |- (m.west); \draw[tree] (0.3,-0.7) |- (l.west); \draw[tree] (0.3,-1.3) |- (r.west);
  \draw[tree] (0.3,-1.9) |- (o.west);
  \draw[tree] (0.9,-2.15) |- (r1.west); \draw[tree] (0.9,-3.35) |- (o1.west);
\end{tikzpicture}
\end{center}

A \hl{record} holds one file's name, size, permissions, dates and a secret
random \hl{seed}. It's padded to 4096 bytes before encryption, so every
record is the same size whatever the name. The file's contents are cut into
\hl{objects} of just under 4\,MiB each. Every object starts with its own ID,
the last one is padded with random bytes up to the next power of two (at
least 64\,KiB), and each is encrypted separately. The objects' names are
worked out from the seed, so only someone who can open the record can tell
which objects belong together.

\shot{on-disk}{Random names, identical sizes, and all ``modified'' at the dawn of Unix time}

\section{Why big files are cut up}

With one object per file, a 300\,MB video would sit in the folder as a
single 512\,MiB blob. Everyone could see there's one large file, roughly how
large, and, since filler is never that big, that it can't be filler. Cut
into pieces, the same video is 72 objects of exactly the same size as the
largest filler and as every other big file's pieces. Nothing on disk says
which pieces belong together, or where one file ends and the next begins.

It's not magic, and some padding and fluff is still doing the work:

\begin{itemize}
  \item \textbf{Small files are still bucketed.} A file under 4\,MiB is one
        object padded to 64\,KiB, 128\,KiB, 256\,KiB and so on, exactly as
        before. The last piece of a big file is padded the same way, so no
        file wastes more than 4\,MiB, where a single padded blob could
        waste nearly half its size.
  \item \textbf{Filler is still the cover story.} The random entries
        \cmd{init} writes are the reason ``objects I can't open'' is
        unremarkable. By default there are only a few tens of megabytes of
        them, enough to hide a few files, not gigabytes. \cmd{vault init
        --filler SIZE} writes more, shaped like a real collection of
        files, at the cost of that much disk.
  \item \textbf{Totals still show.} The number of objects times 4\,MiB is
        a fair estimate of how much you store.
  \item \textbf{Timing still shows.} Someone who watches the folder change,
        or reads a sync service's upload log, sees 72 objects appear together
        and can guess they're one file. The pieces hide a file's size from a
        snapshot, not from someone watching you add it.
\end{itemize}

\section{What an attacker learns}

\begin{vtable}{Y>{\centering\arraybackslash}p{22mm}}
\thead{SOMEONE WITH A COPY OF YOUR VAULT FOLDER\ldots} & \thead{CAN THEY?} \\ \midrule
Read any file's contents & \textcolor{vgreendeep}{\faIcon{times}\ no} \\
Learn file names or folder names & \textcolor{vgreendeep}{\faIcon{times}\ no} \\
Learn exact file sizes & \textcolor{vgreendeep}{\faIcon{times}\ no} \\
Learn when a file was created or changed & \textcolor{vgreendeep}{\faIcon{times}\ no} \\
Tell real entries from decoy or filler & \textcolor{vgreendeep}{\faIcon{times}\ no} \\
Alter a file without it being detected & \textcolor{vgreendeep}{\faIcon{times}\ no} \\
Swap one file's contents for another's & \textcolor{vgreendeep}{\faIcon{times}\ no} \\
Tell which pieces make up one file, from a single copy & \textcolor{vgreendeep}{\faIcon{times}\ no} \\
Count files, and see roughly how much you store & \textcolor{vamberdeep}{\faIcon{check}\ yes} \\
See that the folder changed between two copies & \textcolor{vamberdeep}{\faIcon{check}\ yes} \\
Delete entries & \textcolor{vcoraldeep}{\faIcon{check}\ yes} \\
Put back an older copy of an entry (roll back) & \textcolor{vcoraldeep}{\faIcon{check}\ yes} \\
\end{vtable}

\begin{headsup}
\textbf{Deletion and rollback aren't detectable.} vault can prove a file is
exactly what you stored, but not that it's the \emph{latest} thing you stored.
If the storage isn't under your control, keep an independent backup.
\end{headsup}

\section{What happens during \texttt{add}}

\begin{steps}
  \item Pick a random seed. Read your file 4\,MiB at a time and encrypt each
        piece, the last with random padding, to a hidden temporary file in
        its \file{objects/} folder. Several pieces are written at once.
  \item Give each piece the name the seed dictates. They're now safely
        stored, but nothing refers to them.
  \item Encrypt a record holding the seed, the same way, into \file{records/}.
  \item Only then, delete any older version of the same name.
\end{steps}

Interrupt it anywhere and the worst case is some orphaned encrypted objects
that nothing points to. You never end up with a record pointing at nothing, or
without either version.

\begin{squirrel}
vault's crypto core comes from \emph{Fury}, a larger post-quantum encryption
toolkit. Peek inside any entry and the first line still reads
\cmd{fury-encryption.org/v1}. That's the format header, kept on purpose:
changing working, well-tested cryptographic formats for aesthetic reasons
is how bugs get in.
\end{squirrel}

\begin{pocket}
\begin{pocketlist}
  \item Hybrid post-quantum key wrapping, per-file keys, authenticated streaming encryption.
  \item Fixed-size records, 4\,MiB pieces, power-of-two tails, 1970 timestamps: metadata stays private.
  \item Pieces hide a file's size in any one copy of the folder, not from someone watching it change.
  \item Tampering is caught. Deletion and rollback are not, so keep backups.
\end{pocketlist}
\end{pocket}

% ======================================================================
\chapter{When things go sideways}
\tagline{every message vault can show you, and what to do about it}

\begin{tldr}
Almost every error is one of four things: wrong passphrase, wrong \cmd{--key}
or \cmd{--vault}, a name typo, or a file already in the way. Check the table.
Your data is still fine.
\end{tldr}

\shot{errors}{The usual suspects}

\section{The big table}

\begin{vtable}{>{\ttfamily\footnotesize}p{58mm}Y}
\thead{MESSAGE} & \thead{WHAT IT MEANS, AND THE FIX} \\ \midrule
incorrect passphrase & Typo, Caps Lock, or the wrong key file. Check \cmd{--key} and \env{VAULT\_KEY}. \\
not a vault key file & \cmd{--key} points at something that isn't a vault key, or the key file is damaged. Use your backup. \\
open \ldots/key: no such file or directory & No key there yet. Run \cmd{vault genkey}, or point \env{VAULT\_KEY} at yours. \\
\ldots\ is not a vault; run vault init & Wrong \cmd{--vault} or \env{VAULT\_DIR}, or you haven't run \cmd{init} yet. \\
\ldots\ is already a vault & You ran \cmd{init} on an existing vault. Nothing was changed. \\
NAME: not in the vault & Check the exact name with \cmd{vault list}: it's case-sensitive and includes the folder. Or you typed the other passphrase. \\
\ldots\ already exists & \cmd{get}: the output file exists, so use \cmd{-o}. \cmd{mv}: the new name is taken. \cmd{genkey}: you already have a key. \\
invalid vault path "\ldots" & Names can't start with \cmd{/} or \cmd{../}, end in \cmd{/}, or contain backslashes or control characters. \\
\ldots: not a regular file of at most 1\,TiB & You tried to add a device, a socket or a gigantic file. \\
\ldots\ changed while being added & The file was being written while vault read it. Try again once it's finished. \\
skipping \ldots: not a regular file & Adding a folder: a symlink or special file was skipped. Everything else went in. \\
passphrases do not match & The two entries differed. Nothing was created; run it again. \\
passphrase must not be empty & Blank passphrases aren't allowed, for good reason. \\
decoy passphrase must differ from the passphrase & The decoy can't be the same as the real one. \\
VAULT\_PASSPHRASE is set but empty & Unset it, or give it a value. \\
could not read passphrase (set VAULT\_PASSPHRASE\ldots) & No terminal to ask on, typically cron or CI. Set the variable. \\
waiting for another vault process to finish & Not an error. Another command is writing; this one continues when it's done. \\
vault still locked after a minute\ldots & Plan~9 and WebAssembly only. If nothing else is running, delete the file it names. \\
skipping unreadable record \ldots & A catalogue entry is damaged. Everything else still works; restore \file{records/} from backup to recover that file. \\
object does not match its catalogue record & The encrypted contents were swapped or damaged. Nothing was written. Restore \file{objects/} from backup. \\
object size does not match its catalogue record & Same as above: damage or tampering, caught before any output appeared. \\
cannot find home directory; set VAULT\_KEY & Unusual environments only. Set \env{VAULT\_KEY} and \env{VAULT\_DIR}. \\
unknown command / wrong number of arguments & A typo in the command itself. The usage text follows. \\
\end{vtable}

\section{``My files have vanished!''}

Deep breath. Nine times out of ten, one of these:

\begin{checks}
  \item \textbf{You typed the other passphrase.} The decoy and real sides are
        separate worlds. Try the other one.
  \item \textbf{You're looking at a different vault.} Check \cmd{--vault} and
        \env{VAULT\_DIR}. Is this the vault you think it is?
  \item \textbf{You're looking in a folder.} \cmd{vault list letters} only
        shows what's under \file{letters/}.
\end{checks}

\section{Interrupted in the middle of something?}

Power cut, closed laptop, \keycap{Ctrl}+\keycap{C}. It's fine. vault writes
new data first and removes old data last, so an interruption leaves the vault
consistent. At worst there's an orphaned encrypted object, or a temporary
file starting \file{.pending-}. Both are harmless and both are safe to delete.

\begin{danger}
\textbf{The one thing that can't be fixed:} a lost key file or a forgotten
passphrase. Nothing in this chapter helps with that. Back up your key.
\end{danger}

\begin{pocket}
\begin{pocketlist}
  \item Most errors mean a wrong passphrase, key, vault or name.
  \item ``Missing'' files usually mean the other passphrase or another vault.
  \item Interruptions are safe. Lost keys are not.
\end{pocketlist}
\end{pocket}

% ======================================================================
\chapter{Cheat sheet}
\tagline{the whole tool at a glance. print it, pin it up}

\begin{tcbraster}[raster columns=2,raster equal height=rows,raster column skip=8pt,raster row skip=8pt,
  enhanced,colback=vmist,colframe=vline,boxrule=0.4pt,arc=4pt,left=9pt,right=9pt,top=7pt,bottom=7pt]
\begin{tcolorbox}{\ttfamily\color{vtealdeep}\bfseries vault genkey}\\[2pt]\small Make a key. Asks for a passphrase and an optional decoy.\end{tcolorbox}
\begin{tcolorbox}{\ttfamily\color{vtealdeep}\bfseries vault init}\\[2pt]\small Make an empty vault, with filler. \texttt{--filler 2G} for more.\end{tcolorbox}
\begin{tcolorbox}{\ttfamily\color{vtealdeep}\bfseries vault add FILE|DIR [FOLDER]}\\[2pt]\small Encrypt in, keeping the name. Replaces an existing name.\end{tcolorbox}
\begin{tcolorbox}{\ttfamily\color{vtealdeep}\bfseries vault list [FOLDER]}\\[2pt]\small Size, modified time and name, sorted.\end{tcolorbox}
\begin{tcolorbox}{\ttfamily\color{vtealdeep}\bfseries vault get NAME [-o OUT]}\\[2pt]\small Decrypt out. Never overwrites. Verifies first.\end{tcolorbox}
\begin{tcolorbox}{\ttfamily\color{vtealdeep}\bfseries vault mv OLD NEW}\\[2pt]\small Rename, using the full new name. Instant.\end{tcolorbox}
\begin{tcolorbox}{\ttfamily\color{vtealdeep}\bfseries vault rm NAME}\\[2pt]\small Delete one file. No undo.\end{tcolorbox}
\begin{tcolorbox}{\ttfamily\color{vtealdeep}\bfseries vault help}\\[2pt]\small Show the command summary.\end{tcolorbox}
\end{tcbraster}

\section*{Options and variables}
\begin{vtable}{>{\ttfamily\small}p{40mm}>{\leavevmode\ttfamily\small\color{vvioletdeep}}p{48mm}Y}
\thead{FLAG} & \thead{VARIABLE} & \thead{DEFAULT} \\ \midrule
--key PATH & VAULT\_KEY & \file{\tld/.vault/key} \\
--vault DIR & VAULT\_DIR & \file{\tld/.vault/store} \\
\textrm{(prompt)} & VAULT\_PASSPHRASE & asked on the terminal \\
\textrm{(genkey prompt)} & VAULT\_DECOY\_PASSPHRASE & asked; empty means none \\
-o OUT & \textrm{---} & the file's own name, here \\
--filler SIZE & \textrm{---} & a few tens of MB, for \cmd{init} \\
\end{vtable}

\section*{Numbers worth knowing}
\begin{vtable}{Y>{\raggedleft\arraybackslash}p{40mm}}
\thead{THING} & \thead{VALUE} \\ \midrule
Largest file & 1\,TiB \\
Longest stored name & 1024 bytes \\
Smallest object on disk & 64\,KiB \\
Largest object on disk & 4\,MiB \\
Folders under \file{records/} and \file{objects/} & 256 each \\
Filler entries in a new vault & 8 to 32, or more with \cmd{--filler} \\
Key file size & 319 bytes \\
Exit codes & 0 ok, 1 error, 2 no command \\
\end{vtable}

% ======================================================================
\appendix
\vappendixtrue

\chapter{Glossary}

\begin{description}[style=nextline,leftmargin=0pt,font=\displaybold\color{vink},itemsep=4pt]
  \item[Argon2id] A deliberately slow, memory-hard function that turns your passphrase into a key, making guessing expensive.
  \item[Catalogue] The encrypted list of what's in your vault, stored as one record per file.
  \item[Decoy passphrase] A second passphrase that opens a separate, harmless set of files.
  \item[Entry] One record plus its objects. Real, decoy and filler entries all look identical on disk.
  \item[Filler] Entries sealed to a key that was immediately destroyed. Nobody can open them. They make the decoy plausible.
  \item[FrodoKEM, ML-KEM] Post-quantum key-encapsulation algorithms. vault uses both at once.
  \item[Key file] The small file holding your identity, encrypted under your passphrase (and decoy).
  \item[Object] One encrypted piece of a file, 4\,MiB, or padded to a smaller power of two if it's the last piece.
  \item[Post-quantum] Designed to stay secure against attackers with large quantum computers.
  \item[Record] The encrypted catalogue entry for one file: its name, size, dates and the seed its objects' names come from.
  \item[Vault] The folder where your encrypted files live.
\end{description}

\chapter{About this guide}

This guide lives in \file{docs/guide} and builds itself.

\begin{shell}
$ ./build-pdf.sh    # makes vault-guide.pdf (downloads Tectonic if needed)
$ ./shots.py        # re-shoots every screenshot from the real vault binary
\end{shell}

Every screenshot is real output from the \cmd{vault} program built from the
same source code, run in a throwaway home folder. The only touch-ups: the home
folder is shown as \file{/home/sam}, some lines are highlighted in amber
so they're easier to spot, and the two ``waiting for your passphrase'' shots
are frozen at that moment.

% ======================================================================
\backmatter
\clearpage
\thispagestyle{empty}
\begin{tikzpicture}[remember picture,overlay]
  \shade[top color=vink,bottom color=vnavy] (current page.north west) rectangle (current page.south east);
  \vaultdoor{($(current page.center)+(0,30mm)$)}{38}
  \node[font=\displayblack\fontsize{40}{40}\selectfont,text=white] at ([yshift=-38mm]current page.center) {That's it.};
  \node[font=\lightfont\fontsize{15}{20}\selectfont,text=white!80!vteal,align=center] at ([yshift=-56mm]current page.center)
    {Go encrypt something.\\ (And back up your key.)};
\end{tikzpicture}

\end{document}
